\scanpage{187}
% archive sheet number: 186

La \unsure{catégorie|donnée} d'un \deleted{\unsure{tel}} groupe formel ou groupe pro-algébrique affine \inserted{\unsure{définis comme certains groupes} (ind ($k$-affines))} \unsure{définissent} donc une \deleted{coalgèbre} hyperalgèbre topologique sur $k$, dont l'espace sous-jacent est linéairement compact. Il n'est pas exclu que les hyperalgèbres topologiques ayant cette propriété \deleted{\unsure{puissent}} puissent être interprétées ainsi. \deleted{\illegible{2}}
% slanted insertion at the left, linked to the word affine
% slanted
\margin{\deleted{\unsure{on peut vérifier} \unsure{que} \illegible{4} \unsure{la} \illegible{2}}}

Le mécanisme de dualité de Cartier est alors immédiat, ainsi que son interprétation en termes d'accouplement dans $G_m$. La définition de $F$ et $V$ est aussi immédiate, en car. $p$.

% Sweeping cancellation strokes cross the lower comparison.
\diagram{\begin{tikzcd}[ampersand replacement=\&,column sep=large,row sep=small]
\deleted{G_m} \& \deleted{\mathbb Z/\ell} \& \mu_p \\
\deleted{\unsure{F=x^p}} \& \& \\
\deleted{F(x)=x^p} \& \deleted{F(x)=x^{\deleted{\unsure{p}}}} \& \deleted{F=0} \\
\deleted{V(x)=x} \& \deleted{V(x)=x^p} \& V=\mathrm{id} \\
\& \text{($F$ et $V$ \unsure{absolus|disons})} \&
\end{tikzcd}}
% slanted
\inserted{\unsure{par la réduction naturelle au schéma} : $\mathbb F_p$ \unsure{elle identifie} $G_m=G_m^{(p)}$}
